\section{Limitations and open directions}\label{sec:limits}
\paragraph{Limitations.} These are search results, not optimality proofs, and tightening is a local method. In the main comparison (C07) each instance had 32 runs per method at a fixed budget, and C16 had 16; the $n=12$ cube record came from one of 192 runs, and ``only one method reached it'' holds within the run count and budget stated for that campaign. Gradient paths are matched in steps, not pair evaluations or wall-clock time (hardening was the fastest per run on polygon families); the baselines are matched to a hardening run's pair evaluations, which perturbation--compression overshoots (Table~\ref{tab:diag}). Simulated annealing and perturbation--compression are our implementations; their best runs were less often close enough to be tightened, so the tightening policy favours the gradient paths, most for pc (C14 shows the order survives equal tightening on 28 instances); in C16 the other methods' best runs at the hardening-only instances were tightened after the fact and remained short of the best known value. The lowest-point tuning used instances larger than the test set, and in that round every winner sat on a grid edge or corner and, as pre-registered, the grid was not extended. In 3D only the three gradient paths were compared. Catalogue values come from captions and witnesses as of October 2026; truncated values are only known to their printed precision. Not detecting a difference (between mixed and single-path portfolios, between snap or grow-area and rigid starts, or between the pilot and always-rigid) is not evidence of equivalence at these sample sizes.

\paragraph{Open directions.}
\begin{itemize}\itemsep1pt
\item \emph{A rule that predicts the path.} The pilot showed no detected benefit on held-out families that had no aggregate hardening advantage. A useful rule would have to predict, before running, the families where it does; we have one such family and no rule that singles it out in advance.
\item \emph{External baselines.} Our simulated annealing and perturbation--compression are our own implementations. Running a published nester \cite{Gardeyn2025} or solver \cite{Berthold2026} once, untouched, at the same budget would test whether they were given a fair fight.
\item \emph{Softening as a move.} Re-softening an existing packing and rehardening it, as downstream searches do \cite{n13repo,platonicrepo}, against rigid perturbation from the same packings at equal budget.
\item \emph{Other paths.} Paths through other shape families (an ellipse for a rectangle, a rounded octahedron for a cube, superellipses \cite{Jiao2008}), paths that harden pieces one at a time, or schedules that adapt $\tau$ to the current overlap.
\item \emph{Larger $n$, other solids, other solvers.} Most unsolved instances have $n\ge17$. Our controlled study did not include solids other than cubes; downstream work now explores them \cite{platonicrepo} and higher dimensions \cite{tesseractrepo}, without controlled baselines. Hardening could also seed exact models \cite{Berthold2026} or shrink-and-separate nesters \cite{Gardeyn2025}.
\item \emph{Non-convex and mixed pieces.} The rounded-shape construction needs a convex core; non-convex pieces would need a different family of intermediate shapes, for example a convex decomposition hardened part by part.
\end{itemize}
